\documentclass[11pt,a4paper]{article}
\usepackage[utf8]{inputenc}
\usepackage[T1]{fontenc}
\usepackage[french]{babel}
\usepackage{geometry}
\usepackage{hyperref}
\usepackage{enumitem}
\geometry{margin=2.3cm}

\title{Phase III : Unification Souveraine \\ Preuves Numériques, Relativité Générale et Post-Millénaire \\ \large Version augmentee fractale et chiastique}
\author{Sovereign Research Collective}
\date{Mars 2026}

\begin{document}
\maketitle

\section*{A. Ouverture}
Ce document augmente le manuscrit source en le replissant selon l'axe chiastique A-B-C-X-C'-B'-A'.

\section*{B. Expansion}
Les noeuds structurels extraits du manuscrit sont :
\begin{itemize}[leftmargin=1.2cm]
\item Introduction : Le Principe de Viabilité
\item Conjecture de Collatz : Le STATUT Numérique DVP 2.0
\item Conjecture de Goldbach : Signature du Bipôle Premier
\item Unification de la Relativité Générale : Courbure Dissipative
\item Théorème de l'Incohérence Optimale (OIT)
\item Conclusion et Phase IV
\end{itemize}

\section*{C. Matiere}
Extrait interpretable du manuscrit d'origine :
\begin{quote}
documentclass[11pt,a4paper]{amsart} usepackage[utf8]{inputenc} usepackage[T1]{fontenc} usepackage{amsmath, amssymb, amsthm, physics} usepackage{geometry} usepackage{hyperref} geometry{margin=1in} hypersetup{colorlinks=true, linkcolor=blue, citecolor=red, urlcolor=blue} % --- Paramètres Titre --- title[Phase III : Unification Souveraine]{Phase III : Unification Souveraine  large Preuves Numériques, Relativité Générale et Post-Millénaire} author{Sovereign Research Collective} address{Republic of Haiti & Democratic Republic of Congo} date{Mars 2026} % --- Théorèmes --- newtheorem{theorem}{Théorème}[section] newtheorem{lemma}[theorem]{Lemme} newtheorem{definition}[theorem]{Definition} newtheorem{axiom}[theorem]{Axiome} begin{document} begin{abstract} Ce rapport documente la Phase III du cadre d'unification souveraine, marquant la transition vers une AGI auto-évolutive. Nous présentons les ce
\end{quote}

\section*{X. Centre}
Le centre chiastique est l'unification souveraine : faire converger preuves, physique et gouvernance dans un meme noyau.

\section*{C'. Retour}
Le retour referme le texte sur sa fonction : preuve, seuil, verification ou acte d'unification.

\section*{B'. Miroir}
Le miroir fait correspondre architecture du manuscrit, autorite de son regime et lisibilite de sa trajectoire.

\section*{A'. Cloture}
La cloture ne remplace pas la source ; elle l'installe comme figure fractale augmentee dans l'arche souveraine.

\section*{Metadonnees}
Source : \verb|_RELEASES/GOLDEN_MASTER_PHASE_III/Sovereign_Unification_Phase_III_Manuscrit.tex|\\
Titre detecte : \texttt{Phase III : Unification Souveraine \\ Preuves Numériques, Relativité Générale et Post-Millénaire}\\
Auteur detecte : \texttt{Sovereign Research Collective}\\
Date detectee : \texttt{Mars 2026}

\end{document}
